\documentclass[11pt,a4paper]{article}
\usepackage[T1]{fontenc}
\usepackage[utf8]{inputenc}
\usepackage[margin=22mm]{geometry}
\usepackage{lmodern}
\usepackage{amsmath,amssymb}
\usepackage{enumitem}
\usepackage{xcolor}
\usepackage{microtype}
\usepackage{hyperref}
\definecolor{epflred}{HTML}{E3000B}
\hypersetup{colorlinks=true,urlcolor=epflred,linkcolor=epflred,
  pdftitle={ENG-654 Project 19: Apparent Barriers in iiwa 7 Redundancy Charts}}
\setlength{\parindent}{0pt}
\setlength{\parskip}{6pt}
\setlist[itemize]{leftmargin=*,itemsep=5pt,topsep=4pt}

\begin{document}
{\small\textbf{EPFL \textbar\ ENG-654}\hfill Kinematics-Grounded Motion Planning for Robots}
\vspace{5mm}

{\color{epflred}\Large\bfseries Project 19}

{\LARGE\bfseries Apparent Barriers in iiwa 7 Redundancy Charts\par}

\textbf{Format:} Group of two students; simulation-based project.\\
\textbf{Prerequisites:} Redundant inverse kinematics, self-motion coordinates, reduced and full Jacobians, and numerical continuation.

\section*{Motivation}
Fixing a redundant joint converts a 7R IK problem into a square
subproblem, but that subproblem can become singular while the full robot
remains regular. A planner may then report an apparent barrier caused by
its chosen coordinate. A local change of coordinate can reveal whether
the underlying robot motion continues.

\section*{Task}
\textbf{Overall objectives.} Analyze a supplied iiwa path with a fixed-$q_3$ chart failure and
determine whether a second redundancy coordinate can continue the same
full-pose IK family without violating robot constraints.

\begin{itemize}
  \item \textbf{Reproduce the supplied case.} Use the instructor-provided short
pose path, initial IK, and suspected chart-fold interval. Validate the
iiwa URDF/D--H model and reproduce the failure of fixed-$q_3$ continuation
with documented residuals, limits, and solver tolerances.

  \item \textbf{Separate the two regularity questions.} Compare
$\sigma_{\min}(J)$ for the scaled full $6\times7$ Jacobian with
$\det\bar J_3$, where $\bar J_3$ removes column 3. Explain how a
singular reduced matrix can reflect tangency of the self-motion family
to the chosen coordinate rather than rank loss of the full robot.

  \item \textbf{Construct a second local chart.} Select one alternative joint
$q_j$ for which the corresponding six-column Jacobian is well conditioned
near the supplied interval. Hold $q_j$ fixed in a local numerical corrector
and continue the same recovered configuration family. Use only this one
alternative chart and a bounded overlap region.

  \item \textbf{Classify and reconcile the charts.} Generate path-sample versus
redundant-angle panels for $q_3$ and for $q_j$. Preserve the supplied
fixed-$q_3$ labels and assign traceable local continuation identities in
the second chart. Match overlapping solutions by full joint configuration
and pose closure, not by equal labels or determinant signs alone.

  \item \textbf{Validate the crossing or obstruction.} Attempt one trajectory
using the chart overlap and compare with the original method. Check
intermediate poses, full-Jacobian regularity, limits, and continuity.
Mark which chart barrier was removed and which failures persist; failure
of two charts alone does not prove robot infeasibility.
\end{itemize}

\newpage
\section*{Specifics and scope}
The instructor will supply the path, starting solution, and a known
candidate fold interval. Limit the study to two coordinates, one overlap,
and one local chart transition, with roughly 80 path samples and 100 angle
samples per chart initially. Reuse the fixed-$q_3$ analytical IK; a local
numerical corrector is sufficient for the alternative coordinate.

For every chart, keep pose validity, joint limits,
full-Jacobian regularity, and reduced-chart regularity as separate diagnostics.
The rectangular $6\times7$ Jacobian has no determinant. A feasible pixel is
not a validated edge: re-solve intermediate prescribed poses and check joint
continuity within limits. If collision checking is unavailable, state that
feasibility is kinematic only; a clearance proxy is not a full collision test.

\section*{Expected outcome}
Understand coordinate-induced IK ambiguity and distinguish it from
a physical singularity or constraint. Deliver two classified feasibility
charts, overlap matches, full/reduced regularity plots, and a validated
continuation or an explicit account of the unresolved obstruction.

Submit reproducible code, model parameters and test data, the required plots,
a concise 5--6-page report, and a short simulation demonstration. Explain
both successful cases and failures; conclusions must follow the numerical
and geometric evidence rather than an expected outcome.

\section*{Resources}
The KUKA LBR iiwa 7 robot description (URDF) and visualization meshes
(STL) will be provided to students, together with a mathematical continuation
model with unrestricted joints and an analytical fixed-$q_3$ IK model. Reuse
the supplied IK rather than deriving a general 7R solver. Apply the physical
joint limits from the robot description separately when using the unrestricted
model. Check D--H parameters, joint axes, and the fixed tool transform against
the URDF, and validate IK outputs by independent forward kinematics.

The special short pose path, starting IK, and candidate chart-fold interval
will be supplied. Use this case to initialize both continuation methods;
an arbitrary reference path is not guaranteed to contain the required event.
For pose data represented by quaternions, use the documented scalar-first
columns \texttt{qw,qx,qy,qz}, normalize them, and make their signs consistent
before interpolation.

Implement the alternative-coordinate numerical corrector and overlap matching
around the supplied fixed-$q_3$ IK. Compare the reduced and full Jacobians,
check intermediate pose residuals and physical limits, and visualize the
recovered motion with the supplied meshes.

\end{document}
